\documentclass[12pt]{article}
\usepackage[utf8]{inputenc}
\usepackage[T1]{fontenc}
\usepackage{hyperref}
\title{Yir-Yoront Language Resources for K-12 Curriculum Development}
\author{Compiled by Claude Sonnet 5.5 (claude-sonnet-5-5) for Andr\'as Kornai}
\date{}
\begin{document}
\maketitle

\section{Introduction}
Yir-Yoront (ISO 639-3 yiy; Glottocode yiry1245) is a Pama-Nyungan language of Cape York Peninsula, North Queensland, Australia. Some search results list it under the code yyr. This report is based on one web search round and is short because little teaching material turned up. I found no school curriculum, primer or textbook.

\section{Primary Educational Resources}
No primers, textbooks or school programmes for Yir-Yoront appeared in the search results.

\section{Linguistic and Reference Materials}
\begin{itemize}
\item Barry Alpher, \emph{Yir-Yoront Lexicon: Sketch and Dictionary of an Australian Language} (Mouton de Gruyter, 1991). It holds a grammatical sketch and a dictionary. Catalogue record: \url{https://bvbvf1.bib-bvb.de/vufind/Record/DE-604.BV004553057/Details}
\item Alpher's Yir-Yoront pages at ANU, which include amendments and corrections to the lexicon: \url{https://www.anu.edu.au/linguistics/nash/aust/YY/}
\item Bruce Rigsby's review of Alpher's lexicon, listed in Glottolog: \url{https://glottolog.org/resource/reference/id/6084}
\item Overview article: \url{https://en.wikipedia.org/wiki/Yir-Yoront_language}
\end{itemize}

\section{Religious and Community Resources}
None found.

\section{Digital and Audio Resources}
\begin{itemize}
\item The University of Queensland manuscripts library has an entry on the Yir-Yoront language: \url{https://manuscripts.library.uq.edu.au/index.php/yir-yoront-language}
\end{itemize}
The search summary also said that Paradisec holds Yir-Yoront material in the Capell (AC1) and Alpher (BA1) collections. I did not see a direct URL for it, so none is given.

\section{Curriculum Development Context}
The documentation is mainly academic, led by Alpher's grammar and dictionary. Any curriculum would need to be built from these sources together with archival recordings and the community. No information on current community programmes was found.
\end{document}